\chapter{Security and Access Control}
\index{security}\index{workspace roles}\index{item-level permissions}\index{OneLake data access roles}\index{row-level security}\index{least privilege}\index{privilege creep}\index{separation of duties}%

\begin{objectives}
\begin{itemize}
    \item Assign workspace roles---Admin, Member, Contributor, Viewer---by least privilege.
    \item Distinguish item-level from workspace-level permissions and prefer the narrower grant.
    \item Configure OneLake data access roles and SQL row-level security.
    \item Restrict a user to a single region's data through RLS on the SQL analytics endpoint.
    \item Audit a tenant for privilege creep on a schedule.
    \item Design a healthcare data perimeter with strict separation of duties between engineers and analysts.
\end{itemize}
\end{objectives}

\begin{crosscloud}
Layered access control---platform roles, resource permissions, and in-engine row/column security---is structurally identical across clouds; only the names change.

\vspace{2mm}
{\footnotesize
\begin{tabular}{@{}p{2.8cm}p{3.2cm}p{3.2cm}p{3.2cm}p{3.2cm}@{}}
\toprule
\textbf{Layer} & \textbf{Fabric} & \textbf{AWS} & \textbf{GCP} & \textbf{Snowflake} \\
\midrule
Platform container roles & Workspace roles (Admin/Member/Contributor/Viewer) & IAM policies on resources & IAM roles on projects/folders & Account/database roles \\
Item-level grants & Per-item sharing and permissions & Per-resource policies & Per-resource IAM & Object grants \\
Data-plane scoping & OneLake data access roles & Lake Formation & BigQuery authorized views/dataplex & Row/column policies \\
Row-level security & SQL RLS predicates & Lake Formation cell filters & Row access policies & Row access policies \\
Identity & Microsoft Entra ID & IAM / Identity Center & Cloud Identity & IdP via SCIM/OAuth \\
Audit & Workspace/tenant logs, Purview & CloudTrail & Cloud Audit Logs & Access history \\
\bottomrule
\end{tabular}}

\vspace{1mm}
Databricks governs through Unity Catalog's three-level namespace grants; MongoDB Atlas through database roles and field-level encryption. The portable law: \emph{grant at the narrowest scope that lets the work happen, and audit the drift between policy and reality}.
\end{crosscloud}

\begin{keyterms}
\textbf{Workspace role}: the broad collaboration tier---Admin, Member, Contributor, Viewer---applied to a workspace. \quad
\textbf{Item-level permission}: a grant on one item (lakehouse, report) independent of workspace membership. \quad
\textbf{OneLake data access role}: a role restricting which folders or tables within a lakehouse an identity can read. \quad
\textbf{Least privilege}: granting the minimum access the task requires. \quad
\textbf{Privilege creep}: the slow accumulation of broader-than-needed access over time. \quad
\textbf{Separation of duties}: designing roles so no single identity can both alter data and approve its use.
\end{keyterms}

\section{Week 21 Roadmap}
Week~21 opens Part~VI by securing everything Parts I--V built. Monday covers workspace roles. Tuesday contrasts item-level with workspace-level permissions. Wednesday covers OneLake data access roles and SQL RLS. Thursday is the hands-on project: a workspace with a Contributor role and RLS restricting a user to the US region. Friday designs the healthcare perimeter with separation of duties.

Security in Fabric is layered, and the layers compose: workspace roles set the collaboration ceiling, item permissions narrow it, OneLake roles scope the data plane, and SQL RLS filters the rows \cite{microsoftFabricSecurity}. Misunderstanding any layer leaves a hole the others do not cover.

\begin{notebox}
Production rule for the week: grant access at the narrowest scope---prefer item-level permissions and OneLake data access roles over workspace-wide Member/Admin assignments. Audit quarterly: Viewer on a workspace combined with broad item sharing is the most common form of privilege creep in Fabric tenants.
\end{notebox}

\section{Monday: Workspace Roles}
\index{workspace roles}
Four roles define workspace participation \cite{microsoftFabricWorkspaces,microsoftFabricSecurity}:

\begin{table}[H]
\centering
\small
\begin{tabular}{@{}p{2.6cm}p{5.8cm}p{5.4cm}@{}}
\toprule
\textbf{Role} & \textbf{Can} & \textbf{Give to} \\
\midrule
Admin & Everything: settings, membership, deletion & Workspace owners only; two or three per workspace \\
Member & Create/edit items, manage most settings, add members (if allowed) & Core team leads \\
Contributor & Create and edit items; cannot manage membership or settings & Working data engineers and analysts who build \\
Viewer & Read and consume items & Consumers, stakeholders, auditors \\
\bottomrule
\end{tabular}
\caption{Workspace roles and their intended holders.}
\label{tab:chapter21-roles}
\end{table}

The distinguishing exam-and-practice question is Contributor versus Viewer: a Contributor \emph{changes} items---edits pipelines, modifies notebooks, writes through them---while a Viewer consumes outputs without altering definitions. Most people who ask for Member actually need Contributor; most who ask for Contributor actually need Viewer.

\begin{pitfall}
Admin is not ``senior user.'' An Admin can delete the workspace and everything in it, and can change who else has access. Every unnecessary Admin is an unmonitored disaster path; keep the count at two or three and review it quarterly.
\end{pitfall}

\section{Tuesday: Item-Level versus Workspace-Level Permissions}
\index{item-level permissions}
Workspace roles grant broad capability across every item in the workspace. Item-level permissions grant access to one item---a specific lakehouse, report, or warehouse---to an identity that may have no workspace role at all \cite{microsoftFabricSecurity}.

The decision rule: workspace roles for the team that \emph{builds} the workspace's products; item-level grants for everyone who \emph{consumes} a specific product. A finance analyst who needs the sales report gets the report item---not Viewer on the workspace, which would expose every draft, experiment, and internal artifact the workspace contains.

\begin{tipbox}
When a stakeholder requests ``access to the workspace,'' ask which item they actually need. Item-level grants keep the workspace's internal surface---work in progress, staging tables, experiments---private by default.
\end{tipbox}

Item sharing composes with the deeper layers: sharing a lakehouse item grants the connection, while OneLake roles (Wednesday) scope which data inside it the identity can actually read.

\section{Wednesday: OneLake Data Access Roles and SQL RLS}
\index{OneLake data access roles}\index{row-level security}
\subsection{OneLake Data Access Roles}
OneLake data access roles restrict access to folders or tables \emph{within} a lakehouse: an identity can reach the lakehouse yet read only \texttt{Tables/gold/} and not \texttt{Files/raw/}, or only the virtual folder for their region \cite{microsoftFabricSecurity}. This is the data-plane layer---it applies regardless of which engine (Spark, SQL endpoint, OneLake APIs) the identity uses, which is precisely why it matters: SQL-only security can be bypassed by a Spark notebook unless the storage layer also enforces the rule.

\subsection{SQL Row-Level Security}
Chapter~6 built the mechanism; this week operationalizes it. RLS on the SQL analytics endpoint or warehouse filters rows per user through a predicate function and security policy---the \texttt{US}-region-only analyst is the canonical exercise \cite{microsoftFabricSqlEndpoint,microsoftFabricSecurity}.

\begin{lstlisting}[language=SQL,caption={RLS restricting a principal to US-region rows (Chapter 6 pattern, production framing).},label={lst:chapter21-rls}]
-- Mapping row: the analyst's identity to their region
INSERT sec.user_region (user_principal, region)
VALUES ('analyst-us@contoso.com', 'US');

-- Policy already filters gold.v_sales_enriched via sec.fn_region_predicate.
-- Verify as the restricted principal:
EXECUTE AS USER = 'analyst-us@contoso.com';
SELECT DISTINCT region FROM gold.v_sales_enriched;   -- must return only 'US'
REVERT;
\end{lstlisting}

\begin{pitfall}
Layer-confusion is the recurring security bug: RLS configured on the SQL endpoint does not constrain a Spark notebook reading the same Delta files; a OneLake role does not filter rows inside a permitted table. Each engine's access path needs its layer---test every path a user actually has, not just the one you intended them to use.
\end{pitfall}

\section{Thursday: Hands-on Project}
\index{hands-on lab}\index{row-level security!hands-on}
The Week~21 lab creates a workspace, assigns a Contributor role, and configures RLS on a SQL analytics endpoint so a test user sees only \texttt{US}-region data.

\subsection{Goal and Architecture}
Create \texttt{fabric-de-week21} with a lakehouse containing a regional sales table (reuse any Part~I/II estate). Create or use a test identity for the restricted analyst. Deliverables: the role assignment evidence, the RLS objects, and the two-persona test.

\subsection{Step-by-Step Actions}
\begin{enumerate}
    \item Create the workspace; add the builder identity as Contributor; add the test analyst with item-level access only to the lakehouse.
    \item Create \texttt{sec.user\_region}, the predicate function, and the security policy per Chapter~6.
    \item Insert the analyst's mapping row for \texttt{US}.
    \item As the analyst (or \texttt{EXECUTE AS USER}), query the protected view: only \texttt{US} rows return.
    \item Attempt direct access to the base table: denied.
    \item Attempt the same data via a Spark notebook as the analyst: record what the OneLake layer does and does not block; add a data access role if the scenario demands it.
    \item Document the access matrix: identity $\times$ path $\times$ result.
\end{enumerate}

\subsection{Validation Checklist}
\begin{itemize}
    \item The builder holds Contributor, not Member or Admin.
    \item The analyst holds item-level access, not a workspace role.
    \item RLS filters to \texttt{US} under the analyst identity in SQL.
    \item Base-table access is denied at the object level.
    \item The Spark/OneLake path is tested and its result documented.
    \item The access matrix is written down and reviewed.
\end{itemize}

\section{Friday: Use Case and Architecture}
\index{separation of duties!healthcare}\index{data perimeter}
A healthcare provider must enforce strict separation of duties: data engineers operate pipelines and infrastructure; data analysts query curated data; neither can silently do the other's job; patient data is visible only as far as each role requires.

\subsection{The Perimeter Design}
\begin{itemize}
    \item \textbf{Workspace split}: engineering workspace (Bronze/Silver, pipelines) versus analytics workspace (Gold, semantic models, reports). Engineers are Contributors in engineering and Viewers---at most---in analytics. Analysts have item-level access to curated Gold items only.
    \item \textbf{Data-plane scoping}: OneLake data access roles confine engineers' operational access to the layers they maintain; raw PII folders are readable only by the ingestion identity and the privacy officer.
    \item \textbf{Row and column security}: analysts reach patient-adjacent data through RLS-filtered views with denied sensitive columns (Chapter~6's column-level pattern).
    \item \textbf{Break-glass}: an emergency-access procedure that is time-boxed, dual-approved, and audit-logged---separation of duties includes a governed exception path, because ungoverned ones emerge anyway.
    \item \textbf{Audit}: quarterly privilege-creep review; every grant attributable to a ticket.
\end{itemize}

\begin{pitfall}
The most common separation-of-duties failure is the shared service identity with broad rights used by both teams' automation. Service principals get the same least-privilege treatment as humans---one per pipeline family, scoped to its actual read/write set.
\end{pitfall}

\section{Operational Guidance for Week 21}
Access reviews are operations, not annual theater. Quarterly: list every workspace Admin and Member, every item share, every OneLake role, and reconcile each against a ticket or an owner statement.

\subsection{Performance and Reliability}
RLS predicates add query cost; keep predicate functions simple and indexed where the engine allows. Prefer security by composition (narrow views over RLS'd bases) to baroque per-object grants that nobody can audit.

\subsection{Security and Governance Checklist}
\begin{itemize}
    \item Admins per workspace: two or three, reviewed quarterly.
    \item Consumers get item-level grants; builders get Contributor; Member is exceptional.
    \item OneLake roles cover the Spark path wherever SQL RLS covers the SQL path.
    \item Every RLS policy is tested as a real restricted principal after any change.
    \item Service principals are scoped per workload and vaulted.
    \item Privilege-creep audit is scheduled, evidenced, and acted on.
\end{itemize}

\section{Interview Q\&A}
\begin{enumerate}
    \item \textbf{Q: What distinguishes Contributor from Viewer workspace roles?}\\
    A: Contributors create and modify items---pipelines, notebooks, models; Viewers consume outputs without altering definitions. The distinction is change authority, not seniority.
    \item \textbf{Q: Why prefer item-level permissions over workspace-level grants?}\\
    A: A workspace grant exposes every item---drafts, staging, experiments---to the grantee; item-level grants expose exactly the product the consumer needs and nothing else.
    \item \textbf{Q: What do OneLake data access roles control?}\\
    A: Which folders and tables within a lakehouse an identity can read at the storage plane---enforced regardless of engine, closing the Spark-path bypass of SQL-only security.
    \item \textbf{Q: How does SQL RLS restrict a user to US-region rows only?}\\
    A: A mapping table links the user principal to their region; an inline predicate function evaluates it per row; a security policy binds the predicate to the table---so identical queries return different rows per user.
    \item \textbf{Q: What is privilege creep, and how do you audit for it?}\\
    A: The accumulation of broader-than-needed grants over time---classically Viewer-on-workspace plus broad item sharing. Audit quarterly: enumerate grants, map each to a ticket or owner, and revoke what cannot be justified.
    \item \textbf{Q: Why do service identities need least privilege too?}\\
    A: Automation identities run unattended at scale; a broad shared service principal is an unaudited, always-on privilege that both teams' pipelines inherit---one scoped principal per workload family contains that.
\end{enumerate}

\section{Further Reading}
Review Fabric security fundamentals \cite{microsoftFabricSecurity}, workspace roles \cite{microsoftFabricWorkspaces}, and the SQL endpoint's security model \cite{microsoftFabricSqlEndpoint}. Governance and compliance continue in Chapter~22 \cite{microsoftPurviewFabric}; capacity and deployment controls in Chapter~23 \cite{microsoftFabricGitIntegration}.

\section*{Key Takeaways}
\begin{itemize}
    \item Workspace roles set the collaboration ceiling: Admin is rare, Member is exceptional, Contributor builds, Viewer consumes.
    \item Item-level grants serve consumers; workspace roles serve builders.
    \item OneLake data access roles enforce the data plane across every engine---SQL-only security has a Spark-shaped hole.
    \item RLS composes with object-level security: grant the object, filter the rows, test as the real principal.
    \item Every access path a user actually has must be tested, not just the intended one.
    \item Privilege creep is audited quarterly or it is not audited at all.
    \item Separation of duties includes service identities and a governed break-glass path.
\end{itemize}

\section{Exercises}
\begin{enumerate}
    \item \textbf{(Conceptual)} Construct a scenario where Viewer-on-workspace leaks information that item-level sharing would not.
    \item \textbf{(Conceptual)} Explain why the OneLake layer is necessary even with perfect SQL RLS.
    \item \textbf{(Hands-on)} Extend the lab with a second region and a manager identity mapped to both; test the union behavior.
    \item \textbf{(Hands-on)} Build the privilege-creep audit query/export for your workspace: roles, shares, and last review date.
    \item \textbf{(Design)} Write the break-glass procedure for the healthcare provider: approval, time-boxing, logging, revocation.
    \item \textbf{(Design)} Design the service-principal map for the Part III pipeline estate: one principal per what, scoped to what?
    \item \textbf{(Interview)} In five sentences, explain layered security in Fabric and where each layer's blind spot lies.
    \item \textbf{(Operational)} Write the quarterly access-review procedure with evidence requirements.
\end{enumerate}

\section{Capstone Project: Layered Access Model with Audit Evidence}\label{sec:ch21-capstone}
\index{capstone project}\index{least privilege!capstone}\index{separation of duties!capstone}

\begin{capstonebox}
\textbf{Mission.} Implement the healthcare-grade access model on a working estate: split workspaces, least-privilege roles, item-level consumer grants, OneLake data access roles, tested RLS and column denies, a scoped service-principal map, and a completed quarterly-style audit with evidence.

\textbf{Estimated time.} 5--7 hours.

\textbf{What you'll practice.}
\begin{itemize}
    \item Role assignment by least privilege with an explicit access matrix.
    \item Composing OneLake, object, column, and row security layers.
    \item Testing every engine path per persona.
    \item Producing audit evidence a compliance officer can accept.
\end{itemize}
\end{capstonebox}

\subsection*{Scenario}
Contoso Health Analytics must pass a security review: engineers and analysts share one Fabric tenant but must be provably separated, with patient-identifiable columns denied to analytics and every access decision traceable to a ticket.

\subsection*{Requirements}
\begin{itemize}
    \item Engineering and analytics workspaces with role assignments per the access matrix.
    \item Item-level consumer grants; zero consumer workspace roles.
    \item OneLake data access roles scoping engineers to their layers.
    \item RLS plus a column deny on the patient-adjacent table, tested as each persona.
    \item The service-principal map for all automation.
    \item The audit export: every grant, its justification, and the review sign-off.
\end{itemize}

\subsection*{Implementation Steps}
\begin{enumerate}
    \item Write the access matrix (personas $\times$ assets $\times$ layers).
    \item Assign roles and item grants; configure OneLake roles.
    \item Implement and test RLS and column denies per persona and per engine path.
    \item Map and re-scope service principals.
    \item Run the audit export; reconcile every grant; obtain sign-off.
\end{enumerate}

\subsection*{Deliverables}
\begin{itemize}
    \item The access matrix and role/share inventories.
    \item Test evidence per persona per engine path.
    \item The service-principal map.
    \item The signed audit export.
\end{itemize}

\subsection*{Acceptance Criteria}
\begin{itemize}
    \item[\(\square\)] No consumer holds any workspace role; no engineer holds analytics-workspace write access.
    \item[\(\square\)] RLS and column denies hold under every tested engine path, including Spark.
    \item[\(\square\)] Every service principal is scoped to its workload with vaulted credentials.
    \item[\(\square\)] Every grant in the audit export has a justification or is revoked.
    \item[\(\square\)] Admin counts are two or three per workspace with named owners.
    \item[\(\square\)] No credentials or PII appear in any submitted evidence.
\end{itemize}

\subsection*{Stretch Goals}
\begin{itemize}
    \item Automate the audit export as a scheduled notebook with drift detection against the last review.
    \item Implement the break-glass procedure with time-boxed group membership.
    \item Add an alerting rule on new Admin assignments.
    \item Extend the model with a third persona (external auditor) read-only across both workspaces.
\end{itemize}
